\paragraph{The normalized degeneration and its fields of definition.}
On the complementary period lattice let \(H_B\) be the Hermitian
form, normalized so that its imaginary part is the integral Riemann
form. Let \(C_y\) be the abelian component of the period belonging
to \(y\in Y\). The toric periods are affine in the arc coordinate
\(w\), with linear term given by the positive form of the preceding
construction. The constant term in their pairing can be chosen to be
\begin{equation}\label{eq:slice-cm-normal-pairing}
 T(y,y')=-\frac{H_B(C_y,C_{y'})}{2i}
                 +\operatorname{Im}H_B(\alpha C_y,C_{y'}),
 \qquad \alpha+\bar\alpha=1,
\end{equation}
where \(\alpha\) is the marked integral element of the CM host.
The Rosati relation and the displayed equation for \(\alpha\)
make \(T\) symmetric. They also give

\[
  T(y,\alpha y')+T(y',\alpha y)=T(y,y').
\]

Consequently \(\exp(2\pi i T(y,\alpha y))\) is the even
trivialization with polarization given by \(T\).
Replacing the old symmetric constant pairing by \(T\) is a
translation along the complex cusp flow. Indeed the difference
is conjugate balanced over \(\mathcal O_E\) on \(Y\), so over
the trace pairing it is a complex linear combination of real
parameters of that flow. The slope and abelian projection are
preserved. At large imaginary \(w\) the positive slope absorbs
the fixed translation in each embedding. A Fourier--Jacobi
coefficient under that translation receives its trace character,
which is never zero.

In the normalized model \(n'C_y\) belongs
to the abelian period lattice. Thus
\begin{equation}\label{eq:slice-cm-tau-normal}
  \exp\bigl(2\pi i T(n'y,n'y')\bigr)
       \exp\bigl(\pi H_B(n'C_y,n'C_{y'})\bigr)=1.
\end{equation}
Indeed its remaining exponent divided by \(2\pi i\) is the
integral Riemann pairing of \(\alpha n'C_y\) and \(n'C_{y'}\).
On this sublattice the two gluing points are the identity of the
abelian variety. The exponential product is the comparison of
the polarization of the trivialization with the rigidification
there. This is the canonical torsion normalization.
\label{cor:slice-cm-arc-model-1}

\begin{proposition}[Comparison of the two evaluations on the normalized CM arc]
\label{prop:marked-cm-arc-base-change}
On the marked CM degeneration with pairing
\eqref{eq:slice-cm-normal-pairing}, let \(s,a_\nu\) and \(s',a'_\nu\) denote the positive slope and the Fourier--Jacobi coefficient before and after the change of constant period. Then
\[
 s'=s>0,\qquad a'_\nu=\gamma_\nu a_\nu,\quad\gamma_\nu\ne0.
\]
 Over the fields with the prescribed marking, the
complex and finite-place evaluations use the same toric parameter and
the same normalized theta and frame lines, up to the indicated constant
frame changes.
\end{proposition}

\begin{proof}
\proofstep{Realizing and descending the marked period.}
The marked CM period realization applies to this complementary
lattice with its principal alternating form, split CM type and
bounded basis. It gives a polarized abelian scheme \(B\) after a
finite number field extension and an identification of its complex
periods and CM action with the above ones. In this construction one
keeps the isometry of marked period lattices, so the same
identification carries the Riemann matrix and the Weil pairings on
the gluing points.

Toroidal uniformization for the isotropic
quotient uses \(B\), the character lattice \(Y\), the gluing
homomorphism \(c:Y\to B[n']\), the positive period form and a
trivialization in the pulled-back cubical line. The norm and the
trivialization are compared using the rigidification at zero and
the symmetric primitive theta characteristic. Writing \(E_B\) for the integral period pairing, the symmetric characteristic evaluates on a period \(a\) as
\[
  (-1)^{E_B(\alpha a,a)}.
\]
 The equality of
the period pairings gives equality of these signs, which is needed
for the two Appell--Humbert readings to agree.
\cite{faltings-chai}\label{cor:slice-cm-arc-model-2}

\proofstep{Preserving the toric parameter in descent.}
To descend the normalized model, first adjoin the \(n'\)-torsion
of \(B\). The characteristic-zero torsion scheme is finite, so
all points of \(c\) are now defined over a number field \(k\).
Choose algebraic generators for the one-dimensional cubical
fibres at \(c(y)\). Compare them with the normalized analytic
generators by nonzero scalars \(q(y)\). In the quotient \(\mathbb C^\times/k^\times\), divide the
analytic and algebraic symmetry and cube identities. They
give an even function \(\bar q\) with bimultiplicative
polarization
\[
 b(y,y')=\frac{\bar q(y+y')}{\bar q(y)\bar q(y')}.
\]
The normalized product in \eqref{eq:slice-cm-tau-normal} is
the rigidification over \(k\) on \(n'\)-multiples, so
\[
 b(y,y')^{(n')^2}=b(n'y,n'y')=1.
\]
Bimultiplicativity also gives \(\bar q(0)=1\).
Since \(b(y,-y)=\bar q(y)^{-2}\), one has

\[
  \bar q(y)^{2(n')^2}=1.
\]
 Thus \(q(y)\) satisfies a
power equation over \(k\) and is algebraic.
For a finite integral generating family of \(Y\), it suffices to
adjoin the scalars on the generators and on their pairwise sums.
In the quotient by the enlarged field all polarizations vanish
by bimultiplicativity, and then all the scalars vanish in the
quotient group by the character law. This is one finite extension
for the entire trivialization. The remaining level data are points
in finite torsion schemes of the bounded level and can be adjoined
in a further finite extension.
If \(k_i/k\) are the finitely many fields of definition of the residual level records, their compositum is finite since
\[
  [k_1\cdots k_t:k]\le\prod_{i=1}^t[k_i:k].
\]
 The rigidified line, the marked
action and the normalized toric parameter are transported together.
\label{cor:slice-cm-arc-model-3}

\paragraph{Finite-place realization of the same arc.}

The archimedean construction has now specified the arc and its field. Its evaluation must also extend to the completed model at every finite place needed for the height.

Use a finite place \(w\) of the field \(K'\) just constructed,
write \(V\) for its complete DVR and \(F_V\) for its fraction
field. Base change the abelian part \(B\) to \(V\). By properness
every gluing point in \(B(F_V)\) extends uniquely to \(B(V)\).
The extended points form a homomorphism \(c_V:Y\to B(V)\):
both sides of 
\[
  c_V(y+y')=c_V(y)+c_V(y')
\]
 extend the same
generic point, as do the sides of the CM equivariance identity.
For the rigidified cubical line \(\mathcal T\) on \(B/V\),
the datum retains each generic generator in
\[
             (c_V(y)^*\mathcal T)\otimes_V F_V .
\]

In particular one retains its valuation relative to
\(c_V(y)^*\mathcal T\), even if the generator is not a basis
over \(V\). This is
the datum used in the integral toroidal model. In particular it
has the same generic fibre as the original selected arc. The
level points in the degeneration datum may be kept as generic
level records. Extension of a generic level point as a point of
an abelian torsion group alone would not specify its toric record.

\proofstep{Integral frames and tangent coordinates on the chart.}
An integral framed chart has constant tangent letters in each
CM embedding block and a generator of the theta line on the
integral abelian part. The comparison of two such descriptions
is made over \(F_V\) by base change of the same degeneration
datum.
Thus two ordered tangent bases \(e,e'\) on a common CM block over the generic fibre differ by
\[
  e'=eA,\qquad A\in{\rm GL}_n(F_V),
\]
 Each block of algebraic tangent letters is related to the
integral block by a constant matrix. Constancy follows by taking
invariant tangent vectors of the semiabelian extension and
identifying their pullbacks via the common generic group. The
theta generators there span the same one-dimensional space.

Let
\(\gamma\) be the scalar expressing the algebraic section in the
integral generator.
Thus if \(\vartheta_{\rm alg}\) and \(\vartheta_{\rm int}\) are the indicated generators, their comparison and its valuation are
\[
  \vartheta_{\rm alg}=\gamma\,\vartheta_{\rm int},\qquad v_w(\vartheta_{\rm alg}/\vartheta_{\rm int})=v_w(\gamma).
\]
 These identifications determine the frame
comparison over the entire arc. A change of the torsion
normal form multiplies the reading by a root of unity, as the
torsion sections differ by their level character. This leaves the
finite order of \(\gamma\) unchanged.
\cite{faltings-chai}\label{cor:slice-cm-arc-model-4}

The generator of the global theta ideal at \(w\) is a local
coordinate for the theta line of this same polarized model.
Carry the generator and a Lie frame through the polarized
pullback. The integral theta generator differs from the former
by a unit in \(V\), since both generate the same rank-one module
and the isomorphism is rigidified. If \(t_w\) generates the theta ideal in this trivialization,
then 
\[
  \gamma^{-1}=u t_w
\]
 with \(u\in V^\times\).

Passing to the dual
ideal changes the sign:
\[
 v_w(\gamma)=-v_w(t_w)=v_w I(\mathcal T^\vee).
\] The determinant
of the Lie comparison is that of the same model. This identifies
the central factor of the coefficient line used at infinity as
well as at finite places. In the finite comparison below, the
integral toroidal evaluation is made on smaller punctured discs
over valued extensions. The input is the framed chart extending
the normalized degeneration, with the valuation-compatible
extension of the coefficient field. Base change makes the two
evaluations agree after the constant frame change. The bound on
the integral reading uses the integral lattice of sections on
that chart.
\label{cor:slice-cm-arc-model-5}
\end{proof}
